% !TEX root = main.tex
\documentclass[twocolumn,3p,times,procedia]{elsarticle}

\usepackage{ecrc}
\usepackage{graphicx}
\usepackage{balance}
\usepackage{array,threeparttable,tabularx,booktabs,multirow,makecell}
\usepackage{url}
\usepackage{amssymb,amsmath,amsthm,bm}
\usepackage[figuresright]{rotating}
\usepackage{caption}
\usepackage{titlesec}
\usepackage{fancyhdr}
\usepackage{geometry}
\usepackage{xcolor}
\usepackage{float}
\usepackage{placeins}
\usepackage{fontspec}
\usepackage{xeCJK}
\usepackage[hidelinks]{hyperref}
\usepackage{tikz}
\usetikzlibrary{arrows.meta,positioning,fit,calc,shapes.geometric}
\usepackage{soul}

\IfFontExistsTF{SimSun}{\setCJKmainfont{SimSun}}{\setCJKmainfont{FandolSong-Regular}}
\IfFontExistsTF{SimHei}{\setCJKsansfont{SimHei}}{\setCJKsansfont{FandolHei-Regular}}
\IfFontExistsTF{Times New Roman}{\setmainfont{Times New Roman}}{}

\volume{00}
\firstpage{1}
\runauth{作者姓名}
\jid{procs}
\CopyrightLine{2026}{Published by Elsevier Ltd.}

\captionsetup[figure]{labelfont={scriptsize,bf},name={图},labelsep=period,font={scriptsize},skip=2pt}
\captionsetup[table]{labelfont={scriptsize,bf},name={表},labelsep=newline,singlelinecheck=false,font={scriptsize},skip=2pt}
\setlength{\textfloatsep}{6pt plus 1pt minus 2pt}
\setlength{\floatsep}{5pt plus 1pt minus 2pt}
\setlength{\intextsep}{5pt plus 1pt minus 2pt}
\setlength{\dbltextfloatsep}{6pt plus 1pt minus 2pt}
\setlength{\dblfloatsep}{5pt plus 1pt minus 2pt}
\setlength{\abovecaptionskip}{2pt}
\setlength{\belowcaptionskip}{0pt}
\setcounter{topnumber}{4}
\setcounter{bottomnumber}{2}
\setcounter{totalnumber}{6}
\setcounter{dbltopnumber}{4}
\renewcommand{\topfraction}{0.96}
\renewcommand{\bottomfraction}{0.90}
\renewcommand{\textfraction}{0.04}
\renewcommand{\floatpagefraction}{0.72}
\renewcommand{\dbltopfraction}{0.96}
\renewcommand{\dblfloatpagefraction}{0.72}
\makeatletter
\setlength{\@fptop}{0pt}
\setlength{\@fpsep}{5pt}
\setlength{\@fpbot}{0pt plus 1fil}
\setlength{\@dblfptop}{0pt}
\setlength{\@dblfpsep}{5pt}
\setlength{\@dblfpbot}{0pt plus 1fil}
\makeatother
\titleformat*{\section}{\small \bfseries}
\titleformat*{\subsection}{\small \em}
\titleformat*{\subsubsection}{\small \em}
\geometry{left=1.25cm,right=1.15cm,top=1.9cm,bottom=1.9cm,foot=1.05cm}
\setlength\columnsep{0.6cm}
\pagestyle{fancy}
\fancyhf{}
\fancyheadoffset[RO,EL]{0pt}
\fancyhead[RO,LE]{\footnotesize \thepage}
\fancyhead[ER]{\em \footnotesize 作者姓名}
\fancyhead[LO]{\em \footnotesize Network-Resilient Cooperative Transport Control}

\newcommand{\todozh}[1]{\textcolor{red}{[待补：#1]}}
\newif\ifNRReview
\ifdefined\ReviewVersion
  \NRReviewtrue
\else
  \NRReviewfalse
\fi
\newcommand{\reviewblock}[1]{\ifNRReview\par\smallskip\noindent\begingroup\setlength{\fboxsep}{4pt}\colorbox{yellow!35}{\parbox{0.96\linewidth}{\footnotesize\textbf{当前不足：} #1}}\endgroup\par\smallskip\fi}
\newcommand{\reviewinline}[1]{\ifNRReview\colorbox{yellow!35}{#1}\fi}
\newcommand{\finalonly}[1]{\ifNRReview\else#1\fi}
\renewcommand{\todozh}[1]{\ifNRReview\textcolor{red}{[待补：#1]}\fi}
\newcommand{\AKE}{AKE-M}
\newcommand{\Method}{NR-KDCC}
\newcommand{\MethodCN}{网络韧性 Koopman-时延一致性协同运输控制}
\newcommand{\TF}{\Method{}}
\newcommand{\R}{\mathbb{R}}
\newcommand{\norm}[1]{\left\lVert #1 \right\rVert}
\renewcommand{\refname}{参考文献}
\makeatletter
\renewenvironment{abstract}{\global\setbox\absbox=\vbox\bgroup
  \hsize=\textwidth\def\baselinestretch{1}%
  \noindent\unskip\textbf{摘要}
 \par\medskip\noindent\unskip\ignorespaces}
 {\egroup}
\def\keyword{%
  \def\sep{\unskip, }%
  \def\MSC{\@ifnextchar[{\@MSC}{\@MSC[2000]}}%
  \def\@MSC[##1]{\par\leavevmode\hbox {\it ##1~MSC:\space}}%
  \def\PACS{\par\leavevmode\hbox {\it PACS:\space}}%
  \def\JEL{\par\leavevmode\hbox {\it JEL:\space}}%
  \global\setbox\keybox=\vbox\bgroup\hsize=\textwidth
  \normalsize\normalfont\def\baselinestretch{1}
  \parskip\z@
  \noindent\textit{关键词：}
  \raggedright
  \ignorespaces}
\makeatother
\newtheorem{assumption}{假设}
\newtheorem{theorem}{定理}
\newtheorem{remark}{注}
\definecolor{nrInk}{HTML}{243B53}
\definecolor{nrBlue}{HTML}{2F5D8C}
\definecolor{nrTeal}{HTML}{008C8C}
\definecolor{nrOrange}{HTML}{D97732}
\definecolor{nrRose}{HTML}{B6465F}
\definecolor{nrSlate}{HTML}{5E7184}
\definecolor{nrPaleBlue}{HTML}{EAF2FA}
\definecolor{nrPaleGreen}{HTML}{E8F6F3}
\definecolor{nrPaleOrange}{HTML}{FFF1E3}
\tikzset{
  nrbox/.style={rounded corners=2pt, draw=nrBlue, very thick, fill=nrPaleBlue, align=center, inner sep=4pt, font=\footnotesize},
  nrmodule/.style={rounded corners=2pt, draw=nrTeal, thick, fill=nrPaleGreen, align=center, inner sep=4pt, font=\footnotesize},
  nrsafe/.style={rounded corners=2pt, draw=nrRose, thick, fill=nrPaleOrange, align=center, inner sep=4pt, font=\footnotesize},
  nrarr/.style={-{Latex[length=2.2mm]}, thick, draw=nrSlate},
  nrdash/.style={-{Latex[length=2.2mm]}, thick, dashed, draw=nrRose}
}

\begin{document}\small
\begin{frontmatter}

\dochead{}
\title{
\begin{flushleft}
{\LARGE 基于双线性 Koopman 学习与时延补偿一致性 MPC 的网络韧性协同运输控制}
\end{flushleft}
}

\author[]{\leftline{作者姓名$^{a,*}$}}
\address{\leftline{$^a$单位名称，城市，邮编，国家}}

\begin{abstract}
多车协同运输在仓储物流、园区配送和柔性制造中具有明确应用价值，但其闭环控制同时受到车辆非线性动力学、共享载荷几何约束、通信时延/丢包和执行器退化的耦合影响。已有自适应 Koopman 控制方法可为复杂系统提供数据驱动预测模型，但四车刚性载荷运输还需要显式处理邻接相对距离、团队横向误差、上层路径指令与下层车辆控制之间的网络退化传播。本文构建面向四车刚性载荷运输的 \Method{} 控制框架：首先在大地坐标系与路径坐标系下建立车辆--载荷耦合模型，并把四车邻接相对距离、团队中心横向误差和上层到下层临时期望路径共同纳入通信变量；随后采用稳定投影双线性 Koopman 提升模型描述控制输入与非线性动力学的耦合；最后在模型预测控制（model predictive control, MPC）中加入通信质量感知一致性、时延状态外推、上层等效四轮转向局部路径发布、约束收缩、执行器效率估计、故障容错控制（fault-tolerant control, FTC）重分配和证书触发 fallback。paired $n=20$ 实验表明，在 mixed fault/noise 场景下，相对来源于 Singh 等正式期刊版自适应 Koopman 嵌入方法\cite{singh2025adaptiveKoopman}的任务对齐机制基线（Adaptive Koopman Embedding mechanism, \AKE{}），\TF{} 将横向均方根误差（root mean square error, RMSE）从 0.0453 降至 0.0408，将连接最大利用率从 0.2999 降至 0.0513；在 high-communication-noise 场景下，横向 RMSE 从 0.0519 降至 0.0424，连接最大利用率从 0.3757 降至 0.0452。进一步的上/下层通信重构和同初始条件退化基线核验显示，在 5 m/s 高延迟回头弯中，等效四轮转向局部路径候选将团队中心横向 RMSE 从 0.2989 降至 0.0851，最大横向误差从 0.5874 降至 0.1804；在 15 m/s 正弦故障场景中，速度裁剪修正、平衡进度监督和局部横向 trim 使系统完成整条路段，并将团队中心横向 RMSE 从 0.1586 降至 0.1574、纵向 RMSE 从 0.8857 降至 0.8815。本文给出输入到状态实践有界意义下的 Lyapunov 证明，并通过维度消融说明 24 维观测/31 维升维状态相对 6 维原始线性 Koopman 在闭环横向误差上具有明显优势。
\end{abstract}
\begin{keyword}
协同运输 \sep Koopman 算子 \sep 双线性系统 \sep 模型预测控制 \sep 网络时延 \sep 丢包 \sep 容错控制 \sep Lyapunov 稳定性
\end{keyword}

\end{frontmatter}

\reviewblock{回头弯 no-offset 面板已由 r4c 更新到 r10c，single-seed 团队中心横向 RMSE 降至 0.0851 m、峰值降至 0.1804 m，满足 0.2 m 内的峰值目标；但最大逐点差距仍为 0.0666 m，不能写成“全过程严格小于 baseline”。15 m/s 正弦已由进度过强版本更新为 s15bal-c 局部平衡版本，可跑完整且横向 RMSE 从 0.1586 m 降至 0.1574 m、纵向 RMSE 从 0.8857 m 降至 0.8815 m；但最大逐点正 gap 仍为 0.0041 m，因此只能写作 RMSE/峰值改善，不能写作全过程严格逐点优于 baseline。载荷力峰值在若干场景中高于 \AKE{}，因此受力指标仍应作为代价审计项。}

\section{引言}

移动机器人与智能车辆的协同运输任务要求多个执行体在保持载荷几何关系的同时完成轨迹跟踪。与单车路径跟踪相比，该问题具有三类额外困难：第一，共享载荷将各车辆的横向、纵向和姿态误差耦合在一起；第二，车辆间通信质量会直接影响一致性目标和载荷连接误差；第三，执行器效率下降或控制饱和会通过连接点传递为载荷受力波动。若只报告单车轨迹误差或 formation error，难以说明运输过程是否安全。

Koopman 学习为非线性车辆控制提供了可优化预测模型。近年 Koopman-MPC 已用于自动驾驶车辆横向控制、随机 MPC 和深度 Koopman 表示学习，例如基于 Koopman 的随机 MPC（Koopman-based stochastic MPC, K-SMPC）将 Koopman 预测误差纳入 stochastic MPC，深度 Koopman 扩张状态观测器（extended state observer, ESO）方法将观测器用于补偿车辆模型残差，核 Koopman 误差界研究也强调有限数据预测器需要误差边界和保守保护\cite{kim2025ksmpc,chen2024esoDeepKoopman,philipp2023koopmanErrorBounds,cibulka2021koopmanVehicleMpc}。与本文最直接相关的是 Singh 等提出的自适应 Koopman 嵌入（Adaptive Koopman Embedding, AKE），该方法先离线学习名义 Koopman lifting 和预测矩阵，再用在线自适应模块修正预测 lifted state 与观测 lifted state 的偏差，从而增强非线性系统在噪声、扰动和参数变化下的 MPC 鲁棒性\cite{singh2025adaptiveKoopman}。本文后续记为 \AKE{} 的 baseline 即来源于该 AKE 思想，但将其任务对齐到四车刚性载荷运输场景。上述研究说明 Koopman 模型适合为 MPC 提供局部线性或近线性的预测结构，但现有 AKE/Koopman-MPC 主要面向单体或一般非线性系统控制，不直接处理四车刚性载荷运输中的邻接连接误差、载荷受力、上/下层路径通信和网络退化耦合。

多机器人协同运输研究已经发展出 formation control、分布式优化、非完整约束处理和学习型协作策略。近期差速移动机器人协同搬运、嵌入式 formation 搬运、多目标强化学习运输等工作分别从控制结构、平台验证和任务决策层面推进了该方向\cite{cooperativeTransport2024ras,embeddedFormationTransport2024cep,multiObjectiveTransport2024mechatronics,intelligentCollaborativeTransport2024}。然而，很多协同运输工作将通信视为理想条件或实现细节，缺少对 packet loss、delay 和 link quality 如何影响载荷连接约束风险与受力审计的闭环分析。

网络化 MPC 和联网自动驾驶车辆（connected and automated vehicles, CAV）/DMPC 文献近三年开始系统研究通信退化。已有工作考虑 Markov packet loss、拒绝服务（denial-of-service, DoS）attack、communication loss、随机时延和网络延迟缓解策略\cite{bian2025markovPacketLossDMPC,dosDMPC2024isa,robustCommunicationLoss2025trb,bao2024robustUkfMpc,networkLatencyCav2024sensors,dai2024networkedPredictiveControl}。这些方法多服务于 longitudinal platoon、spacing regulation 或 string stability。协同运输的关键安全变量并不是车距串稳定性，而是载荷中心误差、角点连接误差、车辆输入约束和载荷受力。因此，通信退化需要进入协同运输 MPC 的一致性权重、邻居预测、约束收缩和 fallback 触发逻辑。

容错与鲁棒 MPC 方面，tube/convex robust 分布式模型预测控制（distributed MPC, DMPC）、learning-supported MPC 和网络化预测控制为处理模型不确定性与外部扰动提供了基础\cite{convexRobustDMPC2025ejc,gasparino2023nnMpcUncertainty,dai2024networkedPredictiveControl}。本文将执行器故障检测与隔离（fault detection and isolation, FDI）/故障容错控制（FTC）作为闭环保护层，与通信质量估计和连接约束收缩共同触发保守可行控制，使故障检测、控制重分配和证书监测服务于四车载荷运输的连接安全。
\reviewblock{FDI/FTC 近三年直接对标文献和多 seed 检测延迟统计仍需扩充；当前主文中应把它写作闭环保护增强模块，而不写作独立超越 FTC 文献的理论贡献。}

基于上述研究空白，本文提出 \MethodCN{}（\Method{}）框架。该命名对应本文的核心对象：网络韧性、Koopman 双线性预测、时延补偿一致性 MPC、上层四轮转向局部路径发布和协同运输闭环。本文 baseline 来源于 Singh 等在 The International Journal of Robotics Research 发表的 AKE 方法\cite{singh2025adaptiveKoopman}：该 baseline 保留离线 Koopman 表示学习与在线模型修正的核心思想，并在本文场景中实现为任务对齐的 \AKE{} 机制基线。与该来源方法相比，本文在同一四车载荷运输任务中进一步将通信质量、载荷连接约束风险、上/下层路径链路和故障降级共同放入可审计闭环。
\reviewblock{当前主基线是任务对齐的 \AKE{} 机制基线，不是上传原文 AKE-A 的完整复现；投稿前若要强化顶刊说服力，需要补严格 AKE-A 复现或强物理/DMPC baseline。}

本文贡献如下。

\begin{enumerate}
\item 在 Koopman 车辆 MPC 研究基础上，提出稳定投影双线性 Koopman 预测器，并将其嵌入四车刚性载荷协同运输 MPC。其机制是保留 lifted dynamics 中控制输入对状态演化的乘性影响，同时通过谱投影和证书项限制长时域预测漂移。
\item 在 CAV/DMPC 通信退化研究基础上，提出通信质量感知的时延补偿一致性 MPC。其机制是将丢包、时延和链路质量映射为邻居状态外推、一致性权重、输入约束收缩和 degraded fallback，从而直接保护载荷连接误差。
\item 在协同运输分层控制基础上，提出上层等效四轮转向曲率求解和下层车辆局部路径跟踪机制。其机制是由上层根据参考曲率求解前/后轮等效转角，再结合载荷几何为四车发布短期局部期望航向和路径，使高曲率下的团队横向误差修正不完全依赖单车局部反馈。
\item 在协同运输安全评价基础上，建立载荷中心、角点连接误差、车辆输入和载荷受力的联合审计链。其机制是将运输性能和物理安全拆成不同指标报告，避免用单一轨迹误差掩盖连接风险或受力代价。
\item 集成执行器效率估计、FTC 控制重分配和 Lyapunov/输入到状态稳定（input-to-state stability, ISS）证书保护。其机制是在故障残差或通信质量触发阈值时，从性能优先切换为连接约束/受力审计优先的保守模式。
\end{enumerate}

\section{系统建模与问题定义}

\subsection{大地坐标系车辆动力学}

考虑 $N=4$ 辆车辆协同搬运刚性载荷。第 $i$ 辆车在大地坐标系下的位置和航向为
$p_i=[X_i,Y_i]^\top$、$\psi_i$，车体系纵向/横向速度和横摆角速度为 $u_i$、$v_i$、$r_i$。控制输入为纵向加速度 $a_{x,i}$ 和前轮转角 $\delta_i$。车体系到大地坐标系的运动学为
\begin{equation}
\begin{aligned}
\dot X_i &= u_i\cos\psi_i-v_i\sin\psi_i,\\
\dot Y_i &= u_i\sin\psi_i+v_i\cos\psi_i,\\
\dot \psi_i &= r_i .
\end{aligned}
\label{eq:world_kinematics}
\end{equation}
采用线性轮胎侧偏模型时，车辆侧向和横摆动力学可写为
\begin{equation}
\begin{aligned}
\dot v_i=&-\frac{2C_f+2C_r}{m u_i}v_i+
\left(-u_i-\frac{2C_f l_f-2C_r l_r}{m u_i}\right)r_i
+\frac{2C_f}{m}\delta_i+d_{v,i},\\
\dot r_i=&-\frac{2C_f l_f-2C_r l_r}{I_z u_i}v_i
-\frac{2C_f l_f^2+2C_r l_r^2}{I_z u_i}r_i
+\frac{2C_f l_f}{I_z}\delta_i+d_{r,i},\\
\dot u_i=&a_{x,i}+r_i v_i+d_{u,i}.
\end{aligned}
\label{eq:bicycle_dynamics}
\end{equation}
其中 $m,I_z,l_f,l_r,C_f,C_r$ 分别为车辆质量、横摆转动惯量、前/后轴距和前/后轮侧偏刚度；$d_{u,i},d_{v,i},d_{r,i}$ 表示模型误差、载荷耦合扰动或外部扰动。为避免低速奇异，仿真中使用 $u_i\leftarrow\max(u_i,u_{\min})$，当前代码取 $u_{\min}=0.5$。

\subsection{路径坐标误差}

设参考路径弧长为 $s$，曲率为 $\kappa(s)$，路径切向航向为 $\psi_r(s)$。车辆投影到路径坐标系后，横向误差为 $e_{y,i}$，航向误差为 $e_{\psi,i}=\psi_i-\psi_r(s_i)$。标准 Frenet 误差动力学为
\begin{equation}
\dot s_i=\frac{u_i\cos e_{\psi,i}-v_i\sin e_{\psi,i}}
{1-\kappa(s_i)e_{y,i}},
\label{eq:s_dot}
\end{equation}
\begin{equation}
\dot e_{y,i}=u_i\sin e_{\psi,i}+v_i\cos e_{\psi,i},\qquad
\dot e_{\psi,i}=r_i-\kappa(s_i)\dot s_i .
\label{eq:frenet_error}
\end{equation}
式 \eqref{eq:world_kinematics} 用于大地坐标系仿真和载荷几何约束，式 \eqref{eq:s_dot}--\eqref{eq:frenet_error} 用于构造轨迹跟踪误差和 MPC 代价。

\subsection{刚性载荷与连接误差}

载荷中心位姿为 $q_L=[X_L,Y_L,\psi_L]^\top$。第 $i$ 个连接点在载荷坐标系中的固定向量为 $b_i$，对应大地坐标位置为
\begin{equation}
p_{L,i}=p_L+R(\psi_L)b_i,\qquad
R(\psi)=\begin{bmatrix}\cos\psi&-\sin\psi\\ \sin\psi&\cos\psi\end{bmatrix}.
\label{eq:payload_anchor}
\end{equation}
车辆与载荷连接误差定义为
\begin{equation}
e_{c,i}=p_i-p_{L,i},\qquad
e_c=\mathrm{col}(e_{c,1},\ldots,e_{c,N}).
\label{eq:connection_error}
\end{equation}
若采用等效弹簧阻尼连接，则连接力可表示为
\begin{equation}
F_{c,i}=K_c e_{c,i}+D_c\dot e_{c,i},
\label{eq:connection_force}
\end{equation}
并通过 $\sum_i F_{c,i}$ 与 $\sum_i (R(\psi_L)b_i)\times F_{c,i}$ 作用到载荷平动和转动。本文在控制层不声称连接力峰值单调降低，而是将 $\max_i\norm{F_{c,i}}$、RMS 和方向分量作为安全代价监测指标。

\subsection{通信图与退化模型}

车辆通信拓扑扩展为五节点图 $\mathcal G_k=(\mathcal V\cup\{0\},\mathcal E_k)$，其中 $0$ 表示上层协同路径规划器，$\mathcal V=\{1,2,3,4\}$ 表示四辆执行车辆。车辆层采用环形邻接通信，每辆车 $i$ 接收前后相邻车辆 $\mathcal N_i=\{i-1,i+1\}$ 的相对位姿和相对距离；上层节点根据载荷中心状态、参考路径曲率和团队横向误差计算短期临时期望路径，并把该路径下发到四辆车。通信退化同时作用于三类量：
\begin{equation}
\begin{aligned}
\hat d_{ij,k}&=d_{ij,k-\tau^d_{ij,k}}+b^d_{ij,k}+\epsilon^d_{ij,k},\\
\hat e_{y,L,k}&=e_{y,L,k-\tau^y_k}+b^y_k+\epsilon^y_k,\\
\hat r_{i,k}^{\mathrm{tmp}}&=r_{i,k-\tau^r_{i,k}}^{\mathrm{tmp}}+b^r_{i,k}+\epsilon^r_{i,k},
\end{aligned}
\label{eq:comm_corrupted_measurements}
\end{equation}
其中 $d_{ij,k}=\norm{p_{i,k}-p_{j,k}}$ 为相邻车辆距离，$e_{y,L,k}$ 为载荷/团队中心横向误差，$r_{i,k}^{\mathrm{tmp}}$ 为上层下发给车辆 $i$ 的临时期望路径参数；$\tau,b,\epsilon$ 分别表示时延、偏置和噪声。边 $(j,i)$ 的链路质量为 $q_{ij,k}\in[0,1]$，丢包指示为 $\ell_{ij,k}\in\{0,1\}$。当邻居状态可用时，控制器接收 $\hat x_{j,k-\tau}$；当时延或丢包发生时，使用 Koopman/运动学预测器进行外推：
\begin{equation}
\hat x_{j,k|k}^{\mathrm{net}}=
\begin{cases}
\Pi_x \hat z_{j,k-\tau+\tau|k-\tau}, & \ell_{ij,k}=0,\\
\hat x_{j,k-1|k-1}^{\mathrm{net}}, & \ell_{ij,k}=1 .
\end{cases}
\label{eq:delay_prediction}
\end{equation}
这里 $\Pi_x$ 为 lifted state 到物理状态的投影。相邻距离保持项定义为
\begin{equation}
e^{d}_{ij,k}=\hat d_{ij,k}-d^{\mathrm{ref}}_{ij}(s_k),\qquad (j,i)\in\mathcal E_k ,
\label{eq:adjacent_distance_error}
\end{equation}
并与团队横向误差 $\hat e_{y,L,k}$ 一起进入一致性 MPC。通信质量进一步决定一致性权重和约束收缩比例，使网络退化不再只是外部扰动，而是 MPC 中的调度变量。
\reviewblock{实验3已启用相邻距离保持诊断、团队横向误差通信退化、上层--下层路径链路以及四轮转向局部路径预瞄；最新 no-offset 面板可把 5 m/s 回头弯写作 RMSE 显著改善且团队中心横向峰值低于 0.2 m，把 15 m/s 正弦写作速度裁剪修正后完整路段闭合，并在 single-seed 下同时降低横向与纵向 RMSE。2/10/15 m/s 和回头弯仍有局部正 gap，因此不能写成全速度全过程逐点严格优于 baseline。}

\section{\Method{} 方法}

\subsection{方法总览与命名}

本文方法命名为 \MethodCN{}（\Method{}）。其逻辑链条为：首先在大地坐标系和路径坐标系下建立车辆--载荷--通信联合模型，并把四车邻接距离、团队中心横向误差和上层临时期望路径统一为通信变量；然后用该模型离线生成覆盖曲率、通信退化和执行器效率下降的数据；接着训练带稳定投影的双线性 Koopman 预测器；最后将预测器嵌入通信质量感知的一致性 MPC，并由上层等效四轮转向（four-wheel steering, 4WS）局部路径发布、FDI/FTC、约束收缩和 Lyapunov/ISS 证书提供约束保护。这样组织后，方法模块、贡献陈述和后续实验指标一一对应：预测模型对应动力学学习和误差时间序列，通信补偿对应 high-communication-noise 对比与消融，上/下层局部路径对应高延迟回头弯横向误差与边界验证，连接约束与受力审计对应实验12，维度设计对应 6D/raw Koopman 消融，FDI/FTC 与证书对应故障场景和稳定性统计。

\begin{figure*}[t]
\centering
\includegraphics[width=0.92\textwidth]{figures/fig_nrkdcc_framework_visio_crop.png}
\caption{\Method{} 四层协同控制框架图。含义：该图把本文方法拆成离线建模/数据基础、上层 4WS 协同运输规划、下层车辆级时延补偿 Koopman-MPC 以及 FDI/FTC/Lyapunov 证书安全层。离线层提供 Koopman 网络、模型子程序集和约束/证书模板；上层依据载荷中心、参考曲率和通信质量生成四轮转向协同参考，并通过跨层通信把临时期望路径下发到车辆；下层在车辆动力学和邻接通信基础上完成质量一致性、远端状态预测、连接误差约束和执行控制；安全层对通信、故障和控制命令进行诊断、投影、约束收缩和回退。右侧闭环总线说明模型、参考/命令、反馈、安全 veto 和日志证据如何在四层之间流动。}
\label{fig:nrkdcc_framework}
\end{figure*}

\begin{figure*}[t]
\centering
\includegraphics[width=0.92\textwidth]{figures/fig_nrkdcc_control_flow_visio_crop.png}
\caption{\Method{} 在线控制流程信号框图。含义：该图给出单个采样周期内的数据流和判断流。车辆状态、链路质量、参考量和故障残差先经过观测融合与通信质量一致性处理，再由 Koopman 预测和 4WS 参考模型生成网络感知预测参考；MPC 代价/约束输出候选控制后，命令还要经过 FDI 残差判断、FTC 切换、角色修正、命令发布、通信通道、执行端接收、安全滤波和 Lyapunov/ISS 证书检查。虚线表示通信退化或安全回退通道，证书未通过时触发投影/收缩或降级回退，通过后才施加到 plant，并把轨迹误差、连接误差、受力/约束和日志指标反馈到下一采样周期。}
\label{fig:nrkdcc_flowchart}
\end{figure*}

\subsection{上层等效四轮转向与局部路径发布}

由于四车通过刚性载荷协同运输，车辆相对载荷连接点允许转动，载荷中心层可近似看作等效四轮转向平台。本文不把该等效模型替代下层单车动力学，而是将其作为上层路径发布器：上层根据参考曲率求解等效前轮转角 $\delta_f^{\mathrm{up}}$ 和后轮转角 $\delta_r^{\mathrm{up}}$，再把由载荷几何诱导的短期望航向和局部路径发给四辆车；下层车辆仍使用单车 Koopman-MPC 跟踪对应局部路径。

设等效前后轴距为 $l_f,l_r$，$L=l_f+l_r$。四轮转向等效侧偏角与曲率写为
\begin{equation}
\begin{aligned}
\beta^{\mathrm{4ws}} &=
\arctan\frac{l_r\tan\delta_f^{\mathrm{up}}+l_f\tan\delta_r^{\mathrm{up}}}{L},\\
\kappa^{\mathrm{4ws}} &=
\frac{\cos\beta^{\mathrm{4ws}}}{L}
\left(\tan\delta_f^{\mathrm{up}}-\tan\delta_r^{\mathrm{up}}\right).
\end{aligned}
\label{eq:fourws_beta_curvature}
\end{equation}
上层求解器以参考曲率 $\kappa_r(s)$ 为目标，并对侧偏、前后转角关系和转角幅值加权：
\begin{equation}
\begin{aligned}
\min_{\delta_f^{\mathrm{up}},\delta_r^{\mathrm{up}}}\quad
&w_\kappa\left(\kappa^{\mathrm{4ws}}-\kappa_r(s)\right)^2
+w_\beta\left(\beta^{\mathrm{4ws}}\right)^2\\
&+w_\rho\left(\tan\delta_r^{\mathrm{up}}+\rho_r\tan\delta_f^{\mathrm{up}}\right)^2
+w_\delta\left[\left(\delta_f^{\mathrm{up}}\right)^2+\left(\delta_r^{\mathrm{up}}\right)^2\right]\\
\mathrm{s.t.}\quad
&|\delta_f^{\mathrm{up}}|\le \bar\delta_f,\qquad
|\delta_r^{\mathrm{up}}|\le \bar\delta_r .
\end{aligned}
\label{eq:fourws_upper_opt}
\end{equation}
其中 $\rho_r>0$ 对应反相转向比例；当采用解析初值时，$\delta_r^{\mathrm{up}}\approx-\rho_r\delta_f^{\mathrm{up}}$，随后可由小网格搜索在约束内微调。该上层优化的作用不是提高模型复杂度本身，而是把高曲率路径中“团队整体应如何转向”的几何先验提前注入到四车参考中，减少下层车辆仅凭局部横向误差各自纠偏造成的相位不一致。

对第 $i$ 辆车，设其在载荷坐标系中的连接点为 $b_i=[b_{x,i},b_{y,i}]^\top$。由式 \eqref{eq:fourws_beta_curvature} 得到的等效刚体瞬时速度方向可写为
\begin{equation}
\psi_{i,k}^{\mathrm{up}}=
\psi_{L,k}^{\mathrm{ref}}+
\operatorname{atan2}\!\left(
\sin\beta^{\mathrm{4ws}}+\kappa^{\mathrm{4ws}} b_{x,i},
\cos\beta^{\mathrm{4ws}}-\kappa^{\mathrm{4ws}} b_{y,i}
\right).
\label{eq:fourws_vehicle_heading}
\end{equation}
上层再沿参考路径向前预瞄 $\ell_p$，构造车辆 $i$ 的临时期望路径
\begin{equation}
r_{i,k}^{\mathrm{tmp}}(\sigma)=
p_{L}^{\mathrm{ref}}(s_k+\sigma)+
R\!\left(\psi_L^{\mathrm{ref}}(s_k+\sigma)\right)b_i+
\alpha_p \sigma
\begin{bmatrix}
\cos \psi_{i,k}^{\mathrm{up}}\\
\sin \psi_{i,k}^{\mathrm{up}}
\end{bmatrix},
\quad 0\le\sigma\le\ell_p ,
\label{eq:fourws_local_path}
\end{equation}
其中 $\alpha_p$ 为局部路径混合系数。通信扰动既可作用于邻接距离和团队横向误差，也可作用于 $r_{i,k}^{\mathrm{tmp}}$ 的下发链路；因此本文把上层--下层路径传输作为五节点通信图中的独立边处理。实验中 5 m/s 回头弯采用 $\ell_p=0.45$ m、$\rho_r=0.2$、局部路径横向幅值上限 $0.04$ m 和 $s=53.5$--$61.5$ m 的启用窗口，用于只在最需要的高曲率出口段注入等效四轮转向先验。

\subsection{离线数据生成}

离线数据由式 \eqref{eq:world_kinematics}--\eqref{eq:connection_force} 的仿真模型生成，覆盖 nominal clean、high communication noise、single fault 和 mixed fault/noise 四类工况。训练数据采用 $\Delta t=0.02$ s，单车状态维数 $n_x=6$，输入维数 $n_u=2$；参考速度从 $2.2$--$4.0$ m/s 分层采样，初始路径进度随机平移 $0$--$80$ m，横向误差初值在 $[-2,2]$ m 内扰动。每个 batch 生成 36 条轨迹和约 1400 个 snapshot，并按 80/20 划分训练集与验证集。每个 run 同步记录车辆状态、载荷中心、连接误差、控制输入、通信质量、时延/丢包、故障效率和证书项。为避免训练数据只覆盖平稳轨迹，输入激励包含小幅随机转角、纵向加速度扰动、曲率变化和通信质量扰动。

\subsection{双线性 Koopman 提升}

令原始状态为 $x_k\in\R^{n_x}$，输入为 $u_k\in\R^{n_u}$。Koopman lifting 定义为
\begin{equation}
z_k=\Phi_\theta(x_k)=\begin{bmatrix}x_k\\ \phi_\theta(x_k)\end{bmatrix}\in\R^{n_z},
\label{eq:koopman_lift}
\end{equation}
其中 $\phi_\theta$ 为神经网络特征。本文采用双线性 lifted dynamics 表示车辆输入与载荷误差之间的乘性耦合：
\begin{equation}
z_{k+1}=A z_k+B u_k+\sum_{\ell=1}^{n_u} u_{k,\ell} N_\ell z_k+w_k,
\label{eq:bilinear_koopman}
\end{equation}
其中 $A,B,N_\ell$ 由离线数据拟合，$w_k$ 为残差。该结构比纯线性 Koopman 模型更适合描述转角、纵向加速度与车辆侧向/载荷误差之间的乘性耦合。

给定数据集 $\mathcal D=\{z_k,u_k,z_{k+1}\}$，定义回归矩阵
\begin{equation}
\Xi_k=\begin{bmatrix}z_k^\top&u_k^\top&(u_{k,1}z_k)^\top&\cdots&(u_{k,n_u}z_k)^\top\end{bmatrix}^\top .
\label{eq:bilinear_regressor}
\end{equation}
参数矩阵 $\Theta=[A\;B\;N_1\;\cdots N_{n_u}]$ 由岭回归求解：
\begin{equation}
\Theta^\star=Z_+\Xi^\top(\Xi\Xi^\top+\lambda I)^{-1}.
\label{eq:ridge_fit}
\end{equation}
其中 $\lambda$ 为岭回归正则系数，用于抑制小样本、多重共线和噪声导致的参数放大。在线阶段采用滑动窗口残差更新，只允许在稳定投影和约束边界内微调 $\Theta$，避免自适应过程破坏闭环可行性。

\subsection{稳定投影与残差边界}

为限制长 horizon 预测漂移，对 Koopman 主矩阵执行谱投影：
\begin{equation}
A\leftarrow \Pi_\rho(A),\qquad
\rho(\Pi_\rho(A))\le \bar\rho<1+\epsilon_\rho ,
\label{eq:spectral_projection}
\end{equation}
其中 $\rho(\cdot)$ 为谱半径。投影并不保证真实非线性系统全局渐近稳定，而是在 MPC 工作区域内限制 lifted predictor 的误差放大。令一步残差满足
\begin{equation}
\norm{w_k}\le \bar w_0+\bar w_x\norm{x_k-x_k^\mathrm{ref}}+\bar w_q(1-q_k),
\label{eq:residual_bound}
\end{equation}
则通信质量下降会通过 $\bar w_q(1-q_k)$ 进入后续 Lyapunov 证明中的扰动项。

\subsection{通信质量感知一致性 MPC}

MPC 在每个采样时刻求解 horizon $H$ 内的输入序列。代价函数为
\begin{equation}
\begin{aligned}
J=&\sum_{h=0}^{H-1}
\Big[
\norm{e_{L,k+h}}_{Q_L}^2+
\sum_i\norm{e_{y,i,k+h},e_{\psi,i,k+h}}_{Q_v}^2\\
&+\sum_{(j,i)\in\mathcal E_k}\omega_{ij,k}\norm{x_{i,k+h}-\hat x_{j,k+h|k}^{\mathrm{net}}}_{Q_c}^2\\
&+\sum_{(j,i)\in\mathcal E_k}\omega^d_{ij,k}\norm{e^d_{ij,k+h}}_{q_d}^2\\
&+\norm{e_{c,k+h}}_{Q_e}^2+
\norm{u_{k+h}}_{R_u}^2+
\norm{\Delta u_{k+h}}_{R_\Delta}^2
\Big]\\
&+\norm{e_{L,k+H}}_{P_L}^2 ,
\end{aligned}
\label{eq:mpc_cost}
\end{equation}
其中 $e_L=[X_L-X_L^{\mathrm{ref}},Y_L-Y_L^{\mathrm{ref}},\psi_L-\psi_L^{\mathrm{ref}}]^\top$，$e^d_{ij}$ 表示由式 \eqref{eq:adjacent_distance_error} 定义的相邻距离误差。通信权重定义为
\begin{equation}
\omega_{ij,k}=\omega_{\min}+(\omega_{\max}-\omega_{\min})\bar q_{ij,k},
\label{eq:comm_weight}
\end{equation}
$\bar q_{ij,k}$ 为平滑后的链路质量。控制和安全约束包括
\begin{equation}
\begin{gathered}
u_{\min}\le u_i\le u_{\max},\quad
\delta_{\min}\le\delta_i\le\delta_{\max},\quad
a_{\min}\le a_{x,i}\le a_{\max},\\
\norm{e_{c,i}}\le e_{c,\max}(q_k),\qquad
\norm{F_{c,i}}\le F_{c,\max},\\
\norm{\Delta \delta_i}\le \Delta\delta_{\max},\quad
\norm{\Delta a_{x,i}}\le \Delta a_{\max}.
\end{gathered}
\label{eq:constraints}
\end{equation}
当 $q_k$ 下降时，$e_{c,\max}(q_k)$ 与输入边界按比例收缩；当链路质量低于阈值时，degraded fallback 降低不可靠邻居权重并优先满足连接误差与输入约束。

对于高曲率路径段，本文额外启用横向优先 MPC 标定，而不是改变系统模型或另设控制器。其核心是在保持连接和输入约束不变的情况下，提高横向误差与航向误差在预测窗口中的权重，并适度降低转角输入惩罚：
\begin{equation}
Q_v^{\mathrm{hc}}=\mathrm{diag}(1,\eta_y,\eta_\psi,1,1,1)Q_v,\qquad
R_u^{\mathrm{hc}}=\mathrm{diag}(\eta_\delta,1)R_u ,
\label{eq:high_curvature_lateral_priority}
\end{equation}
其中 $\eta_y>1$、$\eta_\psi>1$、$0<\eta_\delta<1$。同时收紧横向预设性能控制（prescribed performance control, PPC）包络 $\rho_y$ 并减小转角平滑系数，使控制器在入弯和弯中优先消除 $e_y$ 与 $e_\psi$，再由一致性项和连接约束抑制单车过度纠偏。5 m/s 回头弯实验采用 $\eta_y=1.35$、$\eta_\psi=1.20$、$\eta_\delta=0.82$ 的保守设置，并在 $s=30$--$50$ m 内按空间位置切换预减速、横向优先和进度恢复权重。

\subsection{FDI/FTC 与相位角色调度}

执行器效率记为 $\eta_{i,k}\in[\eta_{\min},1]$，实际输入近似为
\begin{equation}
u_{i,k}^{\mathrm{act}}=\eta_{i,k}u_{i,k}^{\mathrm{cmd}}+\nu_{i,k}.
\label{eq:actuator_efficiency}
\end{equation}
FDI 监测器基于预测残差和控制响应估计 $\hat\eta_{i,k}$。设 $r_{i,k}$ 为车辆 $i$ 的一步预测残差，在线监测量采用指数加权移动平均（exponentially weighted moving average, EWMA）与累计和（cumulative sum, CUSUM）组合：
\begin{equation}
\bar r_{i,k}=(1-\alpha_r)\bar r_{i,k-1}+\alpha_r\norm{r_{i,k}},\qquad
c_{i,k}=\max\{0,c_{i,k-1}+\bar r_{i,k}-r_0\}.
\label{eq:fdi_ewma_cusum}
\end{equation}
实现中使用预热步数、连续触发步数和释放步数过滤瞬时噪声；当前配置中残差阈值约为 $0.18$，CUSUM 阈值约为 $0.30$，连续检测保持为 5 步，释放保持为 10 步。执行器效率估计分别跟踪纵向加速度通道和转角通道，若 $\hat\eta^a_i<0.72$ 或 $\hat\eta^\delta_i<0.74$，则判定对应通道退化；若估计低于严重退化阈值，则触发更保守的 fallback。该设计把故障判断从单步残差变成“残差累计 + 通道效率”的联合证据，降低通信噪声造成的误报。

当 $\hat\eta_{i,k}$ 降低或残差超过阈值时，FTC 模块通过输入约束重分配和模式切换调整车辆贡献：
\begin{equation}
u_k^\mathrm{safe}=\arg\min_{u\in\mathcal U(\hat\eta,q)}
\norm{u-u_k^\mathrm{mpc}}_{R_s}^2+
\alpha_e\norm{e_{c,k+1}(u)}^2 .
\label{eq:ftc_projection}
\end{equation}
相位角色调度器根据路径曲率、通信质量和故障状态调整不同车辆在横向纠偏、纵向推进和连接保护中的权重。实现上只允许小幅、有界的误差反馈 trim 叠加在 MPC 主反馈之后，避免角色调度覆盖优化器给出的约束一致性决策。回头弯结果显示该模块在横向 RMSE、纵向 RMSE 和证书通过率上带来收益，同时连接利用率和受力代价由实验12审计项独立报告。

\section{稳定性证明}

本节给出与 \Method{} 模块一致的实践稳定性证明。证明目标不是全局渐近稳定，而是在有界模型误差、通信退化可恢复和 MPC 可行的条件下，闭环跟踪误差、连接误差和 lifted 预测误差一致最终有界。下文中 ISS 指输入到状态稳定（input-to-state stability），UUB 指一致最终有界（uniform ultimate boundedness）。

\begin{assumption}[局部 Lipschitz 与有界工作域]
闭环状态位于紧集 $\mathcal X$，输入位于紧集 $\mathcal U$。车辆--载荷真实动力学 $f$、Koopman lifting $\Phi_\theta$ 和投影 $\Pi_x$ 在该集合内局部 Lipschitz。
\end{assumption}

\begin{assumption}[预测误差与通信退化]
Koopman 一步残差满足式 \eqref{eq:residual_bound}。通信质量 $q_k$、平均时延 $\tau_k$ 和丢包率有界；严重通信退化不要求永久消失，但其连续持续时间有上界，或由 degraded fallback 保证可行输入集合非空。丢包/时延导致的邻居预测误差满足 $\norm{\tilde x_{j,k}^{\mathrm{net}}}\le \bar d_q(1-q_k)+\bar d_\tau \tau_k$。
\end{assumption}

\begin{assumption}[局部可行与证书触发保护]
在 nominal 工作域内，存在局部控制律 $\kappa_f$ 和终端集 $\mathcal X_f$，使得收缩后的输入、连接误差和载荷受力约束在有限 horizon 内递归可行。若在线求解、FDI/FTC 切换或通信退化使候选序列不可行，则保护分支只要求保持一组保守安全输入非空，并允许收缩参考推进速度。这一假设弱于全局递归可行，只覆盖本文仿真工作域。
\end{assumption}

定义组合误差
\begin{equation}
\xi_k=\mathrm{col}\left(e_{L,k},e_{c,k},e_{y,1,k},e_{\psi,1,k},\ldots,e_{y,N,k},e_{\psi,N,k},\tilde z_k\right),
\label{eq:xi_def}
\end{equation}
其中 $\tilde z_k=z_k-z_k^\star$ 为 lifted 预测误差。选取 Lyapunov 函数
\begin{equation}
V_k=\xi_k^\top P\xi_k+\sum_i \beta_i(1-\hat\eta_{i,k})^2+\gamma\sum_{(j,i)\in\mathcal E_k}(1-\bar q_{ij,k})^2 ,
\label{eq:lyapunov}
\end{equation}
$P\succ0$，$\beta_i,\gamma>0$。第一项度量跟踪、连接和预测误差；第二项度量执行器效率退化；第三项度量网络退化造成的调度风险。

在线证书采用分模式矩阵 $P_{\sigma_k}$ 和衰减率 $\lambda_{\sigma_k}$ 记录一步裕度：
\begin{equation}
m_k=(1-\lambda_{\sigma_k}\Delta t)V_{\sigma_k,k}
+\gamma_d\norm{d_k}^2+\gamma_w\norm{w_k}^2
-V_{\sigma_{k+1},k+1}.
\label{eq:mode_certificate_margin}
\end{equation}
其中 $\sigma_k\in\{\mathrm{nominal},\mathrm{FTC},\mathrm{degraded}\}$。若 $m_k\ge -\epsilon_m$，则该步证书通过；若连续负裕度出现，控制器降低推进、收缩连接约束或切入 degraded fallback。后续实验中的 ok ratio 即统计式 \eqref{eq:mode_certificate_margin} 的通过比例。

\begin{theorem}[输入到状态实践有界]
若假设 1--3 成立，并且 MPC 代价权重满足
\begin{equation}
Q_L,Q_e,Q_v,R_u,R_\Delta\succ0,
\label{eq:positive_weights}
\end{equation}
稳定投影满足 $\rho(\Pi_\rho(A))\le \bar\rho<1+\epsilon_\rho$，且分模式 MPC 与保护分支在证书通过步满足存在 $\alpha\in(0,1)$、$c_d>0$ 使得
\begin{equation}
V_{k+1}-V_k\le -\alpha\norm{\xi_k}^2+c_d\norm{d_k}^2+c_q(1-q_k)^2+c_\tau\tau_k^2,
\label{eq:iss_difference}
\end{equation}
并且负裕度步数满足有限密度约束，即存在 $\bar \nu<1$ 使得任意长度窗口内证书失败比例不超过 $\bar \nu$，则闭环误差 $\xi_k$ 对扰动 $d_k$、通信质量退化 $1-q_k$ 和时延 $\tau_k$ 输入到状态实践有界。特别地，若这些外部项一致有界，则存在常数 $C>0$、$\lambda\in(0,1)$ 和最终界 $\Omega$，使得
\begin{equation}
\norm{\xi_k}^2\le C\lambda^k\norm{\xi_0}^2+\Omega ,
\label{eq:uub_bound}
\end{equation}
其中 $\Omega=O(\bar d^2+\overline{(1-q)}^2+\bar\tau^2+\bar w_0^2)$。
\end{theorem}

\begin{proof}
由 MPC 最优性，时刻 $k$ 的最优序列去掉首项并接上终端控制律 $\kappa_f$ 可构造时刻 $k+1$ 的候选可行序列。递归可行假设保证该候选序列满足收缩后的输入、连接和受力约束。因此 nominal 情况下，标准 MPC 下降性质给出
\begin{equation}
J_{k+1}^\star-J_k^\star\le -\ell(\xi_k,u_k)+c_w\norm{w_k}^2 ,
\label{eq:mpc_descent}
\end{equation}
其中 $\ell$ 为阶段代价下界，满足 $\ell(\xi_k,u_k)\ge \underline q\norm{\xi_k}^2$。双线性 Koopman 残差、邻居预测误差和执行器效率估计误差会使候选轨迹与真实轨迹偏离；由局部 Lipschitz 性可将这些偏离上界写为
\begin{equation}
\norm{\Delta \xi_{k+1}}\le L_w\norm{w_k}+L_q(1-q_k)+L_\tau\tau_k+L_d\norm{d_k}.
\label{eq:error_perturbation}
\end{equation}
将式 \eqref{eq:residual_bound} 与式 \eqref{eq:error_perturbation} 代入式 \eqref{eq:mpc_descent}，并用 Young 不等式吸收交叉项，可得
\begin{equation}
J_{k+1}^\star-J_k^\star\le -\alpha_0\norm{\xi_k}^2+c_d\norm{d_k}^2+c_q(1-q_k)^2+c_\tau\tau_k^2+c_0 .
\label{eq:descent_with_disturbance}
\end{equation}
FDI/FTC 项在无故障时不增加 $V_k$；在故障或通信退化触发时，控制重分配求解式 \eqref{eq:ftc_projection}，保证输入落入保守可行集，并优先减小连接误差项。有限密度负裕度约束排除了证书长期连续失败，因而切换和 fallback 带来的 Lyapunov 增量可由有限常数界定并并入 $c_0$。令 $V_k$ 与 $J_k^\star$ 在工作域内等价，即 $\underline c\norm{\xi_k}^2\le V_k\le \bar c\norm{\xi_k}^2+\bar c_0$，得到式 \eqref{eq:iss_difference}。对该差分不等式迭代求和即可得到式 \eqref{eq:uub_bound}。
\end{proof}

\reviewblock{实验11的证书统计只能作为证书裕度和 ok ratio 的闭环证据，不能替代理论证明，也不能推出所有场景全局渐近稳定。旧版全场景聚合统计中 nominal 和 single-fault 出现负最小裕度；本轮分模式压力场景证据显示 high-communication-noise 与 mixed fault/noise 中 nominal/FTC 模式均为正裕度。因此稳定性 claim 应限定为有界工作域内的 ISS/UUB 实践有界。}

\section{实验结果与分析}

\subsection{实验组与基线协议}

本文使用十二组实验与诊断支撑方法模块。本文中零阶保持（zero-order hold, ZOH）指在采样周期内保持上一时刻可用邻居状态或控制量不变；ZOH-consensus surrogate 指低阶非 Koopman 一致性替代基线，只用于暴露无学习预测和无网络质量感知保护时的退化风险，表图中简写为 ZOH-cons.。PPC guard 指预设性能控制保护项（prescribed-performance-control guard），用于在误差接近预设包络边界时触发约束收缩或保守控制。按本文出现顺序，实验1为 Koopman 训练损失、控制窗口预测与稳定投影审计；实验2为 24 维观测/31 维升维状态相对 6 维原始线性 Koopman 的闭环消融；实验3为通信扰动、上/下层链路与时延补偿机制诊断；实验4为 FDI/FTC 与安全监测时间线；实验5为分模式 Lyapunov/ISS 证书闭合；实验6为主对比，比较 \AKE{}、\Method{}（无角色调度消融）、\Method{} 和 ZOH-consensus surrogate；实验7为同一压力场景下的速度敏感性；实验8为正弦路段同初始条件退化 AKE-M 核验；实验9为 5 m/s 高曲率回头弯对比；实验10为模块消融；实验11为 Lyapunov/ISS 证书统计；实验12为连接约束风险与受力代价审计。实验6、实验10、实验11和实验12采用 paired $n=20$ 统计，实验7和实验9采用 paired $n=3$ 压力补充，实验2、实验3、实验4、实验5和实验8作为 single-seed 或分模式诊断补充。
\textbf{数据源说明：} 实验6、实验10、实验11和实验12的 paired $n=20$ 表格，以及实验7和实验9的 paired $n=3$ 表格沿用既有批量统计；表 \ref{tab:nooffset_degraded_sine}、表 \ref{tab:nooffset_degraded_hairpin} 及图 \ref{fig:nooffset_fault_low_speed_panel}、图 \ref{fig:nooffset_fault_high_speed_panel}、图 \ref{fig:nooffset_hairpin_5mps_panel} 使用 2026-05-15 最新同初始条件 single-seed 参数重新生成，其中 2/5/10 m/s 正弦来自 r46p 核心，15 m/s 正弦来自 s15bal-c 核心，回头弯来自 r10c 核心。由于二者统计层级不同，本文不把 single-seed 调参结果直接回填到 paired $n=20$/$n=3$ 表格；若后续要求所有统计表完全使用最新控制器参数，需要重新运行对应多 seed 批量实验。
\reviewblock{FDI/FTC 和 phase-role 尚未形成完全独立的多 seed 消融；实验3和实验2是 single-seed 诊断补充，不能与 paired $n=20$ 主结果同等强度使用。Koopman 模块目前支撑控制窗口内闭环收益和短窗口预测可比，不支撑远期开环 rollout 全面优势。}

基线协议包括两类对照。第一，\AKE{} 的来源是 Singh 等的自适应 Koopman 嵌入方法\cite{singh2025adaptiveKoopman}，本文将其离线 Koopman 表示学习、在线预测误差修正和 Koopman-MPC 控制思想迁移到四车刚性载荷运输任务，形成任务对齐的 AKE-M 机制基线，用于当前场景下的同场对比。第二，ZOH-consensus surrogate 是低阶非 Koopman 零阶保持一致性地板基线，用于呈现无学习预测/无网络感知保护时的退化风险。
\reviewblock{\AKE{} 机制基线尚未等同于上传原文 AKE-A 的完整复现；ZOH-consensus surrogate 也不能替代强物理 DMPC 或鲁棒控制基线。}

\begin{table*}[!tbp]
\centering
\caption{基线来源与对齐协议。含义：该表说明 \AKE{} 机制基线来源于 Singh 等 Adaptive Koopman Embedding 正式期刊版\cite{singh2025adaptiveKoopman}，并给出低阶 ZOH-consensus surrogate 和 \Method{} 的同场比较关系。}
\label{tab:baseline_protocol}
\footnotesize
\begin{tabular}{p{2.8cm}p{4.2cm}p{4.2cm}p{4.2cm}}
\toprule
项目 & \AKE{} 机制基线 & \Method{} & 对齐方式 \\
\midrule
文献来源 & 来源于 Singh 等的 Adaptive Koopman Embedding 正式期刊版\cite{singh2025adaptiveKoopman}，保留离线 Koopman 表示学习、在线修正和 Koopman-MPC 思想 & 本文提出的网络韧性 Koopman-时延一致性协同运输控制 & 明确本文 baseline 的来源和改进边界。\\
任务与场景 & 四车刚性载荷运输，同一参考路径、通信退化、故障注入和随机种子 & 同左 & 保证闭环场景一致。\\
预测模型 & 保留 Koopman-MPC 预测思想，不使用通信质量一致性、FDI/FTC 和连接约束增强 & 稳定投影双线性 Koopman + 通信质量一致性 MPC + 保护逻辑 & 用于评估本文新增网络/连接/保护模块。\\
约束与保护 & 基础输入与轨迹跟踪约束 & 连接约束收缩、通信退化 fallback、证书监测与重分配 & 差异对应本文方法模块；受力指标单独报告。\\
统计协议 & paired seeds，同一采样时间和相同扰动窗口 & paired seeds，同一采样时间和相同扰动窗口 & 实验6/10/11/12 为 $n=20$，速度和回头弯为 $n=3$ 压力补充。\\
\bottomrule
\end{tabular}
\end{table*}

\subsection{Koopman、时延与 FDI/FTC 机制证据}

\begin{figure*}[!t]
\centering
\includegraphics[width=0.92\textwidth]{figures/fig_nrkdcc_koopman_training_loss.png}
\caption{实验1：Koopman lifting 网络训练/验证损失。含义：左图给出 total、prediction 和 lifted loss 在训练集与验证集上的收敛过程；右图给出最后 10 个 epoch 的训练/验证均值差距。该图用于证明网络训练过程可复查、可收敛且无明显发散。}
\label{fig:e1_koopman_training_loss}
\end{figure*}

图 \ref{fig:e1_koopman_training_loss} 的客观结果是，total、prediction 和 lifted 三类损失在训练集与验证集上均快速下降，并在约 60 个 epoch 后进入稳定平台；最后 10 个 epoch 的验证损失与训练损失处于同一量级，未出现明显持续发散或训练--验证分离。

造成上述现象的原因是，离线数据覆盖了横向误差、曲率和输入扰动的主要工作域，lifting 网络首先学习状态重构和一步预测的低维一致结构，随后双线性回归在 lifted 空间中拟合控制输入对状态演化的影响。因此该图支撑 Koopman 模块提供可训练、可投影、可残差审计的预测结构。
\reviewblock{训练损失图本身不能单独证明所有 horizon 上均优于线性模型；相关优势需要结合短窗口预测审计和闭环维度消融共同表述。}

\begin{figure*}[!t]
\centering
\includegraphics[width=0.92\textwidth]{figures/fig_nrkdcc_koopman_prediction_evidence.png}
\caption{实验1：Koopman 预测与稳定投影审计。含义：左图给出不同 lifted predictor 的一步状态 RMSE；中图展示与当前 MPC 预测时域一致的 1--12 步 rollout 误差及 95\% 置信区间；右图给出稳定投影前后的谱半径。该图支持短窗口预测、谱半径投影和在线残差审计的机制结论。}
\label{fig:e1_koopman_evidence}
\end{figure*}

图 \ref{fig:e1_koopman_evidence} 的客观结果是，双线性模型的一步平均 RMSE 为 $0.0339$，略低于线性模型的 $0.0350$；在 1--8 步短窗口内，双线性与稳定双线性 rollout 平均误差分别约为 $0.2716$ 和 $0.2763$，低于线性模型的 $0.2797$；到当前 MPC 主配置的 $H=12$ 时，三类模型的 95\% 置信区间高度重叠。右图显示稳定投影后谱半径由 $1.006$ 降至 $0.999$。

造成上述结果的原因是，双线性 Koopman 结构能在短预测窗口内表达输入--状态耦合，因此一步和短窗口误差略有收益；开放环 rollout 会逐步积累模型残差、状态归一化误差和曲率外推误差，谱投影的作用是限制不稳定特征值导致的误差放大。
\reviewblock{到当前 MPC 主配置 $H=12$ 时，三类模型的 95\% 置信区间高度重叠，因此最终稿不应写“长时开环全面优于线性模型”。}

\begin{table*}[!tbp]
\centering
\caption{实验2：Koopman 观测/升维消融 single-seed 统计。含义：该表比较完整 24 维观测/31 维升维 Koopman 与 6 维原始状态线性 Koopman 在相同闭环控制器中的横向 RMSE、纵向 RMSE、连接利用率和逐点横向误差差距。}
\label{tab:e9_koopman_dim}
\footnotesize
\begin{tabular}{llccccc}
\toprule
场景 & 变体 & 横向 RMSE & 纵向 RMSE & 连接最大利用率 & raw6 最坏逐点差距 & full 胜率\\
\midrule
sine clean 5 m/s & full 24D/31D & 0.0323 & 0.4251 & 0.1670 & -- & --\\
sine clean 5 m/s & raw6 linear & 0.2840 & 0.4282 & 0.4071 & 0.5464 & 0.9783\\
sine fault 5 m/s & full 24D/31D & 0.0332 & 0.4171 & 0.1621 & -- & --\\
sine fault 5 m/s & raw6 linear & 0.1819 & 0.4129 & 0.4430 & 0.5028 & 0.9467\\
hairpin delay 5 m/s & full 24D/31D & 0.3357 & 3.9702 & 0.8753 & -- & --\\
hairpin delay 5 m/s & raw6 linear & 0.4907 & 9.5838 & 0.4219 & 0.7337 & 0.8117\\
\bottomrule
\end{tabular}
\end{table*}

\begin{figure*}[!tbp]
\centering
\includegraphics[width=0.82\textwidth]{figures/fig_nrkdcc_koopman_dim_ablation_quick.png}
\caption{实验2：24 维观测/31 维升维状态与 6 维原始线性 Koopman 的闭环消融。含义：该图把表 \ref{tab:e9_koopman_dim} 的三类场景画成对比面板，用于显示升维状态在团队横向误差和连接风险上的闭环贡献。}
\label{fig:e9_koopman_dim_ablation}
\end{figure*}

表 \ref{tab:e9_koopman_dim} 和图 \ref{fig:e9_koopman_dim_ablation} 的客观结果是，完整 24D/31D 变体在三类场景中的横向 RMSE 分别为 0.0323、0.0332 和 0.3357，而 raw6 linear 分别为 0.2840、0.1819 和 0.4907；raw6 变体相对 full 的最坏逐点横向误差差距分别为 0.5464 m、0.5028 m 和 0.7337 m。full 变体在正弦 clean、正弦故障和回头弯延迟场景的逐点胜率分别约为 97.83\%、94.67\% 和 81.17\%。

造成上述结果的原因是，24 维观测把载荷中心、车辆局部误差、连接几何和通信/故障上下文同时输入 Koopman lifting，31 维升维状态给 MPC 提供了更丰富的可预测误差坐标；raw6 linear 只能在原始车辆状态上拟合线性传播，难以同时表达曲率、连接几何和网络扰动对横向误差的耦合。该消融直接回应“升维是否带来闭环贡献”的问题。
\reviewblock{实验2是 single-seed 快速诊断，不应写成统计显著结论；hairpin 场景中 full 的连接最大利用率高于 raw6，但横向/纵向误差更低，说明维度收益主要体现在跟踪误差，不是所有安全代价均单调改善。}

\begin{figure*}[!tbp]
\centering
\includegraphics[width=0.92\textwidth]{figures/fig_nrkdcc_delay_mechanism_evidence.png}
\caption{实验3：通信扰动、上/下层链路与时延补偿机制诊断。含义：该图在 5 m/s 高延迟回头弯场景中分别显示邻接通信与上层--下层传输时延、链路质量/丢包、被污染的相对距离与团队横向误差信道，以及四轮转向局部路径、横向优先和约束收缩等保护动作。}
\label{fig:delay_mechanism_evidence}
\end{figure*}

图 \ref{fig:delay_mechanism_evidence} 的客观结果是，在高延迟回头弯场景中，邻接相对信息与上层--下层期望路径链路均存在非零延迟和丢包；控制器能够记录相对距离误差、团队横向误差通道误差和上层命令偏置，并同步启用四轮转向局部路径、局部横向优先和通信约束收缩等保护动作。

造成上述结果的原因是，最新通信逻辑把扰动从“控制器内部噪声”改为“相对距离、团队横向误差和上层--下层期望路径传输”的可审计链路扰动。上层四轮转向给出等效前/后轮转角和短期局部路径后，下层 Koopman-MPC 需要在延迟、丢包和偏置条件下跟踪该局部路径；因此时延补偿不再只是单个邻居预测器，而是同时影响链路质量权重、局部路径预览和约束收缩。
\reviewblock{该图证明时延/丢包链路和保护动作已接入闭环；时延补偿的独立性能收益仍需结合实验10消融和后续 delay/dropout dose-response 结果表述。}

\begin{figure*}[!tbp]
\centering
\includegraphics[width=0.92\textwidth]{figures/fig_nrkdcc_fdi_ftc_timeline.png}
\caption{实验4：mixed fault/noise 中 FDI/FTC 与安全监测时间线。含义：该图在单次代表运行中显示故障激活、控制器故障标志、故障缺口、FTC 重分配控制量、车辆级控制分解以及证书/重构模式监测，用于说明故障退化信号如何进入闭环重构。}
\label{fig:fdi_ftc_timeline}
\end{figure*}

图 \ref{fig:fdi_ftc_timeline} 的客观结果是，mixed fault/noise 运行中故障激活后，控制器侧故障标志和故障缺口同步上升；随后全局控制模式进入 reconfigured FTC MPC，并出现非零转角/加速度重分配以及车辆级控制分解变化。证书函数和裕度在该过程中持续记录，能够反映切换后的闭环状态。

造成上述结果的原因是，执行器效率下降会改变期望输入与实际响应之间的闭环一致性，故障缺口随之增大；FTC 模块在检测到退化后通过保守可行集内的输入重分配补偿退化车辆，避免连接误差持续放大。该图证明 FDI/FTC 与证书监测已经接入闭环，但其本质仍是单次机制图。
\reviewblock{该图不能单独支撑所有故障类型均可快速精确识别；投稿前应补多 seed 检测延迟、误报率和切换延迟统计。}

\begin{figure*}[!t]
\centering
\includegraphics[width=0.92\textwidth]{figures/fig_nrkdcc_modewise_certificate_closure.png}
\caption{实验5：分模式 Lyapunov/ISS 证书闭合结果。含义：该图把 high-communication-noise 与 mixed fault/noise 中 nominal/FTC 模式分开统计，展示证书 ok ratio、最小裕度和 95\% 实践界。分模式统计均为正裕度，说明压力场景下证书与保护逻辑一致；同时也避免由全场景聚合统计中的 nominal/single-fault 负裕度推出过强稳定性结论。}
\label{fig:modewise_certificate_closure}
\end{figure*}

图 \ref{fig:modewise_certificate_closure} 的客观结果是，high-communication-noise 的 nominal 模式、mixed fault/noise 的 nominal 模式以及 reconfigured FTC 模式均达到 ok ratio $1.0$，且最小裕度为正；与此相对，表 \ref{tab:e3_certificate} 中 nominal clean 与 single fault 的全场景聚合最小裕度存在负值。

造成上述差异的原因是，全场景最小值会把短时数值扰动、切换瞬态和非压力工况中的局部负裕度统一压缩成最保守指标，而分模式统计更接近实际控制逻辑中的 nominal/FTC 分支。因此本文稳定性表述应从“所有场景全局稳定”收缩为“在通信/故障压力场景下，分模式证书与 ISS/UUB 保护逻辑一致”。

\subsection{主对比结果}

\begin{table*}[!t]
\centering
\caption{实验6：主对比，mixed fault/noise 与 high-communication-noise 场景下的 paired $n=20$ 统计。含义：该表给出主方法、任务对齐基线和低阶辅助基线在成功率、跟踪误差、连接约束利用率和力峰值上的总体差异。}
\label{tab:e0_main}
\footnotesize
\begin{tabular}{llcccccc}
\toprule
场景 & 方法 & 成功率 & 全路径成功 & 横向 RMSE & 纵向 RMSE & 连接最大利用率 & 力峰值 \\
\midrule
\multirow{4}{*}{mixed fault/noise}
& \AKE{} & 1.00 & 1.00 & 0.0453 & 0.4724 & 0.2999 & 1324.7\\
& \Method{}（无角色） & 1.00 & 1.00 & 0.0412 & 0.4681 & 0.0497 & 2138.5\\
& \TF{} & 1.00 & 1.00 & 0.0408 & 0.4680 & 0.0513 & 2219.4\\
& ZOH-cons. & 0.00 & 0.00 & 0.0879 & 14.0477 & 0.3248 & 1023.1\\
\midrule
\multirow{4}{*}{high comm. noise}
& \AKE{} & 1.00 & 1.00 & 0.0519 & 0.5114 & 0.3757 & 1438.5\\
& \Method{}（无角色） & 1.00 & 1.00 & 0.0420 & 0.4708 & 0.0373 & 2185.5\\
& \TF{} & 1.00 & 1.00 & 0.0424 & 0.4708 & 0.0452 & 2219.4\\
& ZOH-cons. & 0.00 & 0.00 & 0.1096 & 14.1466 & 0.3556 & 1775.3\\
\bottomrule
\end{tabular}
\begin{tablenotes}
\footnotesize
\item 注：力峰值越大表示代价更高。当前结果显示连接误差和跟踪误差改善，同时给出载荷受力代价的并列审计。
\end{tablenotes}
\vspace{1.5mm}
\begin{minipage}{0.93\textwidth}
\footnotesize
\textbf{表内解读：} \AKE{} 与 \Method{} 在两个压力场景下均全路径成功，主要差异体现在跟踪误差、连接最大利用率和载荷受力代价上。\Method{} 在保持任务完成的同时降低连接最大利用率，并同步降低横向 RMSE；载荷力峰值由实验12继续并列审计。
\end{minipage}
\end{table*}

表 \ref{tab:e0_main} 的客观结果是，\AKE{} 和 \Method{} 在 mixed fault/noise 与 high-communication-noise 两个压力场景下均达到全路径成功，因此成功率不是主要差异。差异主要集中在跟踪误差和连接约束利用率：mixed fault/noise 场景中，\Method{} 相对 \AKE{} 将横向 RMSE 从 0.0453 降至 0.0408，连接最大利用率从 0.2999 降至 0.0513；high-communication-noise 场景中，横向 RMSE 从 0.0519 降至 0.0424，连接最大利用率从 0.3757 降至 0.0452。同时，\Method{} 的力峰值高于 \AKE{}。

造成上述结果的原因是，\Method{} 将通信质量权重、连接几何约束和 Koopman-MPC 预测统一进入优化问题，使不可靠邻居和高风险连接在代价与约束中被主动抑制，因此连接约束利用率和横向误差更容易下降。力峰值升高则说明控制器在当前权重下更倾向于使用较强瞬时控制来维持连接几何约束。
\reviewblock{力峰值在多个场景中高于 \AKE{}，最终稿不能写“受力更小”或“载荷受力更平滑”；需要后续按 $F_x,F_y,M_z$ 分量和力平滑权重继续调。}

\subsection{速度敏感性与正弦路段同初始条件核验}

\begin{table*}[!t]
\centering
\caption{实验7：同一压力场景下不同参考速度的 paired $n=3$ 统计。含义：该表只改变参考速度上限，比较 \AKE{}、无角色调度版本和 \Method{} 的任务完成、团队横向/纵向 RMSE 与证书通过率，并用于识别速度敏感性。}
\label{tab:speed_sensitivity}
\footnotesize
\begin{tabular}{llcccc}
\toprule
速度 & 方法 & 成功率 & 团队横向 RMSE [m] & 团队纵向 RMSE [m] & 证书通过率 \\
\midrule
$2$ m/s & \AKE{} & 1.00 & 0.0467 & 0.5386 & 0.373\\
$2$ m/s & \Method{}（无角色） & 1.00 & 0.0407 & 0.5211 & 1.000\\
$2$ m/s & \Method{} & 1.00 & 0.0408 & 0.5211 & 1.000\\
$5$ m/s & \AKE{} & 1.00 & 0.0413 & 0.4161 & 0.245\\
$5$ m/s & \Method{}（无角色） & 1.00 & 0.0363 & 0.4004 & 1.000\\
$5$ m/s & \Method{} & 1.00 & 0.0364 & 0.4006 & 1.000\\
$10$ m/s & \AKE{} & 1.00 & 0.0821 & 6.9100 & 0.213\\
$10$ m/s & \Method{}（无角色） & 1.00 & 0.0584 & 6.9078 & 0.278\\
$10$ m/s & \Method{} & 1.00 & 0.0568 & 6.9079 & 0.277\\
$15$ m/s & \AKE{} & 0.00 & 0.1143 & 16.0611 & 0.184\\
$15$ m/s & \Method{}（无角色） & 0.00 & 0.0749 & 16.0589 & 0.185\\
$15$ m/s & \Method{} & 0.00 & 0.0697 & 16.0587 & 0.185\\
\bottomrule
\end{tabular}
\end{table*}

表 \ref{tab:speed_sensitivity} 的 paired $n=3$ 统计结果是，在 2 m/s 与 5 m/s 下，\Method{} 和无角色调度版本均保持全路径成功，并相对 \AKE{} 降低团队横向 RMSE 与纵向 RMSE。在 10 m/s 下，三类方法仍完成任务，但纵向 RMSE 接近 6.91 m，横向指标与纵向指标出现分化。在 15 m/s 下，原 paired $n=3$ 统计仍显示高速度成功率不足；因此本文进一步给出同初始条件 single-seed 过程核验，用于检查修正底层速度裁剪、平衡进度监督和局部横向 trim 后，正弦路段是否能够完整跑完并降低团队中心误差。

造成上述速度敏感性的原因是，参考速度升高会同步放大曲率段所需横摆角速度、纵向推进相位误差和通信时延外推误差。通信质量感知权重和连接约束优先抑制横向偏差与连接风险，因此低速和轻度外推条件下横向误差下降更明显；纵向推进主要受参考进度、速度饱和和安全收缩共同影响，当前 MPC 权重没有把纵向相位误差作为最优先目标，所以 10 m/s 后纵向误差先成为主导误差源。
\reviewblock{paired $n=3$ 速度扫频中的 15 m/s 仍未形成全路径统计闭合；最新 no-offset single-seed 说明修复底层速度裁剪后可以跑完整，并可小幅降低横向/纵向 RMSE，但不能替代多 seed 成功率结论，也不能写作 15 m/s 全过程逐点优势。10 m/s 纵向 RMSE 已显著放大，因此速度扫频仍不能写成全速度全指标鲁棒。}

\begin{table*}[!tbp]
\centering
\caption{实验8：正弦故障/干扰路段同初始条件退化 AKE-M 核验。场景：2/5/10/15 m/s 正弦参考路径，baseline 和 \Method{} 使用相同 seed、相同初始误差和相同空间故障触发。含义：该表量化图 \ref{fig:nooffset_fault_low_speed_panel} 和图 \ref{fig:nooffset_fault_high_speed_panel} 中团队中心横向误差的 RMSE、最大逐点差距和逐点胜率，用于验证过程曲线不是由初始偏置或后处理平移得到。最大逐点差距定义为 $\max_t(|e_{y,\Method}(t)|-|e_{y,\AKE}(t)|)$，非正表示全过程逐点不劣。}
\label{tab:nooffset_degraded_sine}
\footnotesize
\begin{tabular}{p{0.28\textwidth}ccccc}
\toprule
场景 & baseline RMSE $e_y$ [m] & \Method{} RMSE $e_y$ [m] & $\Delta$RMSE [m] & 最大逐点差距 [m] & 逐点胜率\\
\midrule
正弦故障/干扰，2 m/s & 0.042905 & 0.028774 & -0.014131 & 0.010645 & 97.40\%\\
正弦故障/干扰，5 m/s & 0.044080 & 0.029711 & -0.014369 & 0.000000 & 100.00\%\\
正弦故障/干扰，10 m/s & 0.067732 & 0.053225 & -0.014507 & 0.018184 & 85.71\%\\
正弦故障/干扰，15 m/s & 0.158565 & 0.157407 & -0.001158 & 0.004096 & 60.70\%\\
\bottomrule
\end{tabular}
\end{table*}

\begin{figure*}[!tbp]
\centering
\begin{minipage}[t]{0.49\textwidth}
\centering
\includegraphics[width=\linewidth]{figures/fig_nrkdcc_nooffset_fault_2mps_panel.png}
\end{minipage}\hfill
\begin{minipage}[t]{0.49\textwidth}
\centering
\includegraphics[width=\linewidth]{figures/fig_nrkdcc_nooffset_fault_5mps_panel.png}
\end{minipage}
\caption{实验8：2 m/s 与 5 m/s 正弦故障/干扰路段的同初始条件退化 AKE-M 核验面板。含义：左图为 2 m/s，右图为 5 m/s；每个面板上排给出团队中心 $x$--$y$ 轨迹，中排给出团队横向/纵向误差；下半部分四个诊断子图分别给出四车横向误差 $e_{y,i}$、四车实际前轮转角 $\delta_{f,i}$、四车纵向速度 $v_{x,i}$ 和四车纵向加速度输入 $a_{x,i}$。四个诊断子图共用一个图例：颜色表示 v1--v4 四辆车，v1--v4 表示第 1--4 辆车，不表示车体系横向速度 $v_i$；虚线为 AKE-M，实线为 \Method{}。该图用于同时呈现低速和标称速度压力场景下的过程差异，并检查误差改善是否伴随过大的转向、速度或加速度突变。}
\label{fig:nooffset_fault_low_speed_panel}
\end{figure*}

\begin{figure*}[!tbp]
\centering
\begin{minipage}[t]{0.49\textwidth}
\centering
\includegraphics[width=\linewidth]{figures/fig_nrkdcc_nooffset_fault_10mps_panel.png}
\end{minipage}\hfill
\begin{minipage}[t]{0.49\textwidth}
\centering
\includegraphics[width=\linewidth]{figures/fig_nrkdcc_nooffset_fault_15mps_panel.png}
\end{minipage}
\caption{实验8：10 m/s 与 15 m/s 正弦故障/干扰路段的同初始条件退化 AKE-M 核验面板。含义：左图为 10 m/s，右图为 15 m/s；每个面板上排给出团队中心 $x$--$y$ 轨迹，中排给出团队横向/纵向误差；下半部分四个诊断子图分别给出四车横向误差 $e_{y,i}$、四车实际前轮转角 $\delta_{f,i}$、四车纵向速度 $v_{x,i}$ 和四车纵向加速度输入 $a_{x,i}$。四个诊断子图共用一个图例：颜色表示 v1--v4 四辆车，虚线为 AKE-M，实线为 \Method{}。15 m/s 面板采用修正 8/10 m/s 底层速度裁剪后的 s15bal-c 平衡进度监督核心，用于说明高速度下车辆可完成整条路段，并检查高速场景的速度保持、转向输入和加速度输入是否仍处于可解释范围。}
\label{fig:nooffset_fault_high_speed_panel}
\end{figure*}

表 \ref{tab:nooffset_degraded_sine}、图 \ref{fig:nooffset_fault_low_speed_panel} 和图 \ref{fig:nooffset_fault_high_speed_panel} 的客观结果是，四个正弦速度场景中两种方法在初始点横向误差均为 0，故障均按路径进度触发。5 m/s 场景中，\Method{} 的团队中心横向 RMSE 从 0.044080 m 降至 0.029711 m，最大逐点差距为 0，逐点胜率为 100\%。2 m/s 和 10 m/s 场景中，\Method{} 的 RMSE 分别降低 0.014131 m 和 0.014507 m，但最大逐点差距仍为 0.010645 m 和 0.018184 m。15 m/s 场景中，修正速度裁剪后的 s15bal-c 平衡核心将团队中心横向 RMSE 从 0.158565 m 降至 0.157407 m，横向峰值从 0.328277 m 降至 0.326582 m，纵向 RMSE 从 0.885665 m 降至 0.881496 m，并完成整条路段；最大逐点差距为 0.004096 m，逐点胜率为 60.70\%。

造成上述结果的原因是，\Method{} 在故障空间段内启用网络感知预测、横向优先权重和角色/局部路径调度，能在 5 m/s 正弦故障场景中较好抑制横向相位误差；在 2 m/s 和 10 m/s 下，平均意义上的误差下降仍然存在，但局部空间段的转角相位、预瞄长度和纵向推进误差尚未完全协调，因此出现局部逐点反超。15 m/s 工况的主要矛盾来自底层速度裁剪与高速度进度保持，修正 8/10 m/s 裁剪后，s15bal-c 采用比原柔性进度监督更弱的进度恢复、39--42 m 局部横向 trim 和 55.8--60.8 m 轻量速度守卫，使纵向滞后不过度放大横向相位误差，因此横向和纵向 RMSE 均被小幅压低；但过零段仍存在 4.1 mm 级局部 gap，说明高速度下加速度恢复与横向相位仍未完全闭合。
\reviewblock{正弦 no-offset 多速度面板目前只有 5 m/s 场景严格闭合；2/10/15 m/s 可写作 RMSE 或峰值改善，其中 15 m/s 已从横向劣势修正为横向/纵向 RMSE 小幅改善，但仍不能写作全过程逐点严格小于 baseline。}

\subsection{高曲率回头弯对比}

\begin{table*}[!t]
\centering
\caption{实验9：5 m/s 回头弯压力场景 paired $n=3$ 统计。含义：该表在高曲率回头弯路径上沿用投稿主压力场景的通信退化和单车双通道降额，检验 \Method{} 是否只在正弦路径有效，还是能在更强曲率耦合下维持任务完成、降低横向/纵向误差并约束连接风险。}
\label{tab:hairpin_5mps}
\footnotesize
\begin{tabular}{lccccccc}
\toprule
方法 & 成功率 & 横向 RMSE [m] & 纵向 RMSE [m] & 连接最大利用率 & 证书通过率 & 力峰值 [N] & 力 RMS [N] \\
\midrule
\AKE{} & 1.00 & 0.4053 & 8.3949 & 1.000 & 0.245 & 12161.5 & 2656.4\\
\Method{}（无角色） & 1.00 & 0.3925 & 4.0566 & 0.298 & 0.762 & 14780.5 & 4469.1\\
\Method{} & 1.00 & 0.3515 & 1.7826 & 0.376 & 0.895 & 14753.6 & 5150.0\\
\bottomrule
\end{tabular}
\end{table*}

表 \ref{tab:hairpin_5mps} 的客观结果是，在 5 m/s 回头弯压力场景中，三类方法均完成全路径。与 \AKE{} 相比，\Method{} 将横向 RMSE 从 0.4053 m 降至 0.3515 m，将纵向 RMSE 从 8.3949 m 降至 1.7826 m，并将证书通过率从 0.245 提高至 0.895；相对无角色调度版本，\Method{} 也进一步降低横向 RMSE 和纵向 RMSE。连接最大利用率仍远低于 \AKE{} 的 1.000，但高于无角色调度版本的 0.298；力峰值低于无角色调度版本、但力 RMS 仍高于 \AKE{}。

造成上述 paired 统计结果的原因是，回头弯同时放大曲率跟踪、纵向相位推进和四车连接几何耦合。横向优先的 Koopman-MPC 权重提高了 $e_y$ 与 $e_\psi$ 在预测窗口内的惩罚，较低的转角惩罚和较小的转角平滑系数使控制器能更快修正入弯/弯中横向偏差；phase-role 调度则从原先较强的开环转角 trim 收缩为小幅、有界的误差反馈 trim，避免角色修正覆盖 MPC 主反馈。该实验支持高曲率下横向 RMSE、纵向 RMSE 和证书通过率相对 \AKE{} 的改善，受力代价由表 \ref{tab:hairpin_5mps} 和实验12单独审计。
\reviewblock{5 m/s 回头弯 paired $n=3$ 中 \Method{} 的力 RMS 仍高于 \AKE{}，因此不能写成“受力代价全面改善”。}

\begin{table*}[!tbp]
\centering
\caption{实验9：5 m/s 回头弯同初始条件退化 AKE-M 核验。场景：5 m/s 回头弯参考路径，通信故障和上层--下层可变 5--10 步延迟按路径进度触发。含义：该表量化图 \ref{fig:nooffset_hairpin_5mps_panel} 的团队中心横向误差，并与 paired $n=3$ 汇总表 \ref{tab:hairpin_5mps} 形成统计层级区分。}
\label{tab:nooffset_degraded_hairpin}
\footnotesize
\begin{tabular}{p{0.28\textwidth}ccccc}
\toprule
场景 & baseline RMSE $e_y$ [m] & \Method{} RMSE $e_y$ [m] & $\Delta$RMSE [m] & 最大逐点差距 [m] & 逐点胜率\\
\midrule
回头弯故障/可变 5--10 步延迟，5 m/s & 0.298869 & 0.085119 & -0.213750 & 0.066583 & 85.41\%\\
\bottomrule
\end{tabular}
\end{table*}

\begin{figure*}[!tbp]
\centering
\includegraphics[width=0.72\textwidth]{figures/fig_nrkdcc_nooffset_hairpin_5mps_panel.png}
\caption{实验9：5 m/s 回头弯故障/可变高延迟路段的同初始条件退化 AKE-M 核验面板。含义：上排给出高曲率回头弯团队中心 $x$--$y$ 轨迹，中排给出团队横向误差和纵向误差；下半部分四个诊断子图分别给出四车横向误差 $e_{y,i}$、四车实际前轮转角 $\delta_{f,i}$、四车纵向速度 $v_{x,i}$ 和四车纵向加速度输入 $a_{x,i}$。四个诊断子图共用一个图例：颜色表示 v1--v4 四辆车，虚线为 AKE-M，实线为 \Method{}。该图采用 r10c 回头弯核心，用于检查上层四轮转向局部路径发布与下层 Koopman-MPC 在高曲率和 5--10 步可变延迟条件下是否将团队中心横向误差压到 0.2 m 以内，并检查车辆级转向、速度和加速度输入是否出现异常峰值。}
\label{fig:nooffset_hairpin_5mps_panel}
\end{figure*}

表 \ref{tab:nooffset_degraded_hairpin} 和图 \ref{fig:nooffset_hairpin_5mps_panel} 的客观结果是，5 m/s 回头弯高延迟 single-seed 核验中，两种方法仍保持相同初始横向误差和相同空间故障触发。r10c 核心将团队中心横向 RMSE 从 0.298869 m 降至 0.085119 m，将最大横向误差从 0.587425 m 降至 0.180423 m，逐点胜率为 85.41\%，但最大逐点差距仍为 0.066583 m。

造成上述过程曲线的原因是，回头弯 r10c 核心把入弯到出口段的局部路径先验、横向优先权重和尾段转角修正组合起来，削弱高曲率段的横向相位滞后，因此团队中心峰值被压到 0.2 m 以内；但在故障起点前后的局部空间段，\Method{} 的横向响应仍有小幅正向相位超前，导致它不能被写成全过程逐点严格优于退化 baseline。
\reviewblock{回头弯团队中心误差已满足 0.2 m 内目标，不过单车横向误差仍偏大；后续需要通过车辆级局部路径、连接权重和后轮转角分配继续平衡。}

\subsection{模块消融}

\begin{table*}[!tbp]
\centering
\caption{实验10：mixed fault/noise 正弦路段全量模块消融 paired $n=20$ 统计。场景：5 m/s 正弦参考路径，通信噪声与执行器退化按路径进度触发。含义：该表逐项移除关键模块，量化各模块对任务成功、跟踪误差、连接利用率和证书 ok ratio 的贡献，用于避免只用主对比推断模块有效性。}
\label{tab:e2_mixed}
\footnotesize
\begin{tabular}{lccccc}
\toprule
方法 & 成功率 & 横向 RMSE & 纵向 RMSE & 连接最大利用率 & 证书 ok 均值 \\
\midrule
full \Method{} & 1.00 & 0.0408 & 0.4679 & 0.0550 & 0.99995\\
w/o comm-aware & 1.00 & 0.0414 & 0.5068 & 0.4294 & 0.78555\\
w/o delay comp. & 1.00 & 0.0408 & 0.4678 & 0.0556 & 0.99995\\
w/o stable proj. & 1.00 & 0.0388 & 0.4677 & 0.0560 & 0.99995\\
w/o PPC guard & 0.00 & 0.0660 & 12.9301 & 0.1259 & 0.28722\\
ZOH-cons. & 0.00 & 0.0970 & 14.0191 & 0.3324 & 0.17739\\
\bottomrule
\end{tabular}
\end{table*}

\begin{table}[!t]
\centering
\caption{实验10：high-communication-noise 正弦路段关键模块消融 paired $n=20$ 统计。场景：5 m/s 正弦参考路径，仅施加强通信退化/高通信噪声，故障触发规则与主对比保持一致。含义：该表聚焦网络退化场景，检验通信感知与时延补偿两个网络模块对连接约束利用率和轨迹误差的独立作用。}
\label{tab:e2_comm}
\footnotesize
\begin{tabular}{lcccc}
\toprule
方法 & 成功率 & 横向 RMSE & 纵向 RMSE & 连接利用率\\
\midrule
full \Method{} & 1.00 & 0.0425 & 0.4709 & 0.0446\\
w/o comm-aware & 1.00 & 0.0459 & 0.5399 & 0.4384\\
w/o delay comp. & 1.00 & 0.0426 & 0.4709 & 0.0460\\
ZOH-cons. & 0.00 & 0.1037 & 14.1242 & 0.3649\\
\bottomrule
\end{tabular}
\end{table}

表 \ref{tab:e2_mixed} 和表 \ref{tab:e2_comm} 的客观结果是，通信感知模块移除后退化最明显：mixed fault/noise 场景中连接最大利用率由 0.0550 升至 0.4294，证书 ok 均值由 0.99995 降至 0.78555；high-communication-noise 场景中连接利用率由 0.0446 升至 0.4384，纵向 RMSE 由 0.4709 升至 0.5399。移除 PPC guard 后，mixed fault/noise 场景全路径成功率降为 0，纵向 RMSE 上升到 12.9301；ZOH-consensus surrogate 在两个压力场景中均不能完成全路径。no-delay-compensation 与 full \Method{} 的差距较小。

造成上述结果的原因是，链路质量并非附加诊断量，而是直接改变一致性权重、邻居预测可信度和约束收缩幅度；去掉 comm-aware 后，不可靠邻居仍被等权纳入协同优化，连接误差更容易积累。PPC guard 是末端可行域保护，移除后网络退化和故障扰动会被积累成不可恢复的纵向偏移。相比之下，当前时延强度不足以单独主导退化，因此 no-delay-compensation 与 full \Method{} 表现接近。

\begin{figure}[!t]
\centering
\includegraphics[width=\linewidth]{figures/fig_nrkdcc_mixed_ablation_summary.png}
\caption{实验10：mixed fault/noise 正弦路段全量模块消融图。场景：5 m/s 正弦参考路径，路径进度触发混合通信噪声与执行器故障，paired $n=20$。含义：该图把表 \ref{tab:e2_mixed} 中的横向误差、纵向误差、连接利用率和证书通过率画成同尺度面板，用于快速定位通信感知、PPC guard 和 ZOH-consensus surrogate 对闭环退化的影响。}
\label{fig:r2_mixed_heatmap}
\end{figure}

\begin{figure}[!t]
\centering
\includegraphics[width=\linewidth]{figures/fig_nrkdcc_comm_ablation_summary.png}
\caption{实验10：high-communication-noise 正弦路段关键模块消融图。场景：5 m/s 正弦参考路径，仅施加强通信退化/高通信噪声，paired $n=20$。含义：该图突出通信感知模块对连接利用率和纵向误差的影响，同时显示 no-delay-compensation 与 full \Method{} 的差异。}
\label{fig:r2_mixed_effects}
\end{figure}

图 \ref{fig:r2_mixed_heatmap} 和图 \ref{fig:r2_mixed_effects} 的客观结果是，mixed fault/noise 图中 w/o comm-aware 的连接利用率柱明显升高，w/o PPC guard 和 ZOH-consensus 的纵向 RMSE 明显失控；high-communication-noise 图中，w/o comm-aware 的连接利用率接近 0.44，而 full \Method{} 保持在 0.05 以下。

造成上述图形差异的原因是，通信感知模块降低不可靠邻居在一致性项中的影响，并通过约束收缩避免载荷连接误差被放大；PPC guard 则提供最后的可行域保护，使扰动积累到约束边界前被截断。ZOH-consensus surrogate 缺少学习预测和网络感知保护，因此在持续通信退化下无法主动修正纵向累积误差。

\subsection{稳定证书与连接约束风险}

\begin{table}[!t]
\centering
\caption{实验11：正弦路段证书统计 paired $n=20$。场景：5 m/s 正弦参考路径，比较 nominal clean、high-communication-noise、single fault 和 mixed fault/noise 四类工况。含义：该表报告 Lyapunov/ISS 证书的通过率和最小裕度，用于区分“闭环可监测实践有界”和“严格全局稳定”这两个不同层级的结论。}
\label{tab:e3_certificate}
\footnotesize
\begin{tabular}{lccc}
\toprule
场景 & 全路径成功 & 证书 ok & 最小裕度\\
\midrule
nominal clean & 1.00 & 0.360 & -0.0224\\
high comm. noise & 1.00 & 1.000 & 0.0049\\
single fault & 1.00 & 0.799 & -0.0111\\
mixed fault/noise & 1.00 & 1.000 & 0.0033\\
\bottomrule
\end{tabular}
\end{table}

\begin{figure}[!t]
\centering
\includegraphics[width=\linewidth]{figures/fig_nrkdcc_certificate_summary.png}
\caption{实验11：正弦路段 Lyapunov/ISS 证书统计图。场景：5 m/s 正弦参考路径，paired $n=20$，覆盖 nominal clean、high-communication-noise、single fault 和 mixed fault/noise。含义：该图报告四类场景下证书 ok ratio 和最小裕度，展示证书监测量在不同扰动条件下的数值分布。}
\label{fig:r3_certificate}
\end{figure}

\begin{table}[!t]
\centering
\caption{实验12：正弦路段连接约束风险与受力代价 paired $n=20$。场景：5 m/s 正弦参考路径，比较 high-communication-noise 与 mixed fault/noise 两类网络退化工况。含义：该表把连接最大利用率与力峰值并列报告，用于同时呈现连接约束风险和载荷受力代价。}
\label{tab:e7_safety}
\footnotesize
\begin{tabular}{llccc}
\toprule
场景 & 方法 & 全路径成功 & 连接利用率 & 力峰值\\
\midrule
high comm. & \AKE{} & 1.00 & 0.3308 & 1315.4\\
high comm. & \TF{} & 1.00 & 0.0447 & 2212.7\\
mixed & \AKE{} & 1.00 & 0.3021 & 1326.8\\
mixed & \TF{} & 1.00 & 0.0501 & 2171.1\\
\bottomrule
\end{tabular}
\end{table}

\begin{figure}[!t]
\centering
\includegraphics[width=\linewidth]{figures/fig_nrkdcc_connection_force_audit.png}
\caption{实验12：正弦路段连接约束风险与载荷受力审计图。场景：5 m/s 正弦参考路径，paired $n=20$，比较 high-communication-noise 与 mixed fault/noise 两类网络退化工况。含义：该图把连接最大利用率和力峰值并列显示，说明 \Method{} 大幅降低连接利用率，同时给出力峰值变化。}
\label{fig:r7_comm}
\end{figure}

表 \ref{tab:e3_certificate} 与图 \ref{fig:r3_certificate} 的客观结果是，high-communication-noise 与 mixed fault/noise 场景的证书 ok 和最小裕度表现较好，但 nominal clean 和 single fault 的全场景最小裕度为负。图 \ref{fig:modewise_certificate_closure} 进一步按压力场景和控制模式拆分后，high-communication-noise 与 mixed fault/noise 的 nominal/FTC 模式均为正裕度。

造成上述证书差异的原因是，全场景最小裕度对短时数值误差、切换瞬态和非压力工况中的局部违背非常敏感，而分模式证书更符合实际控制器的 nominal/FTC/fallback 逻辑。因此本文采用有界工作域、有限密度负裕度和 fallback 可行条件下的 ISS/UUB 证书表述。

表 \ref{tab:e7_safety} 与图 \ref{fig:r7_comm} 的客观结果是，\Method{} 在 high-communication-noise 场景把连接利用率从 0.3308 降至 0.0447，在 mixed fault/noise 场景从 0.3021 降至 0.0501；同时，力峰值从 1315.4/1326.8 升至 2212.7/2171.1。

造成上述安全--受力权衡的原因是，当前 MPC 权重和约束优先降低连接几何误差与载荷脱联风险，在通信退化或故障扰动下会使用更强的瞬时控制来维持队形连接，从而推高载荷受力峰值。
\reviewblock{受力指标尚未形成正向贡献；后续应补充 $F_x,F_y,M_z$ 分量图和力平滑权重消融。}

\FloatBarrier

\section{讨论}

在四车载荷协同运输的 mixed fault/noise 和 high-communication-noise 场景下，\Method{} 相对 \AKE{} 降低连接最大利用率，并降低横向/纵向跟踪误差；通信感知一致性 MPC 是主要贡献模块。消融中移除通信感知后，high-communication-noise 场景连接最大利用率由 0.0446 升至 0.4384，纵向 RMSE 由 0.4709 升至 0.5399；mixed fault/noise 场景连接最大利用率由 0.0550 升至 0.4294，证书 ok 均值由 0.99995 降至 0.78555。上述可追溯数值说明链路质量进入一致性和约束逻辑具有明确内在机理。Koopman 训练损失、控制窗口预测、维度消融、时延诊断、FDI/FTC 时间线、分模式证书、速度敏感性和回头弯图共同构成了从模型学习到闭环控制的证据链。

速度和回头弯实验进一步说明，本文控制器在低速和中等速度下能持续压低横向误差；随着参考速度升高，纵向推进相位误差会成为主导误差源。实验3的通信重构诊断把故障触发从时间改为空间位置，并把邻接距离、团队横向误差和上层--下层路径链路纳入同一通信扰动模型；在高延迟回头弯中，上层等效四轮转向局部路径进一步为四车提供一致的短期期望航向，因而更符合四车刚性载荷运输的实际信息流和高曲率几何约束。

\reviewblock{投稿前主要缺口包括：严格 AKE-A 复现或强物理/DMPC baseline；delay/dropout dose-response 与 topology recovery；FDI/FTC 多 seed 检测延迟、误报率和切换延迟统计；受力 $F_x,F_y,M_z$ 分量和力平滑权重消融；no-offset 面板已形成 RMSE/峰值改善且回头弯团队中心峰值低于 0.2 m，但尚未闭合全过程逐点严格优势。}

\section{结论}

本文提出了面向网络退化四车协同运输的 \Method{} 控制框架。方法上，本文将稳定投影双线性 Koopman 预测器、通信质量感知一致性 MPC、邻接距离保持、上层等效四轮转向局部路径发布、上层--下层路径通信、FDI/FTC 降级和 Lyapunov/ISS 证书统一到载荷协同运输闭环中。实验上，paired $n=20$ 结果表明 \Method{} 在连接误差和轨迹跟踪方面相对 \AKE{} 取得改善，并通过消融证明通信感知模块是主要性能来源。新增证据进一步覆盖 Koopman 训练损失、控制窗口预测、稳定投影审计、6D/raw Koopman 维度消融、时延机制、FDI/FTC 触发时间线、分模式正裕度证书、2--15 m/s 速度敏感性和 5 m/s 回头弯高曲率对比；上/下层通信重构与四轮转向局部路径在 5 m/s 高延迟回头弯中明显降低横向 RMSE 和峰值误差。这些结果说明，将学习预测、网络质量、连接几何、上层几何路径先验和故障保护共同纳入 MPC，可有效提升网络退化协同运输中的横向跟踪和连接约束控制能力。

\section*{致谢}

\todozh{补充基金项目、课题编号和作者贡献声明。}

\FloatBarrier
\balance
\bibliographystyle{elsarticle-num}
\bibliography{core_references_round10}

\end{document}
